\documentclass[11pt]{article}
\usepackage[T1]{fontenc}
\usepackage{lmodern,microtype,amsmath,amssymb,amsthm,mathtools,mathrsfs,tikz-cd}
\usepackage[margin=1in]{geometry}
\usepackage{xurl,hyperref}
\hypersetup{colorlinks=true,linkcolor=blue,citecolor=blue,urlcolor=blue,
pdftitle={Local connectedness, range functors, and essential incompleteness},pdfauthor={Florian Lengyel}}
\newtheorem{theorem}{Theorem}[section]
\newtheorem{proposition}[theorem]{Proposition}
\newtheorem{lemma}[theorem]{Lemma}
\newtheorem{corollary}[theorem]{Corollary}
\newcommand{\N}{\mathbb N}
\newcommand{\PA}{\mathrm{PA}}
\newcommand{\Pfn}{\mathbf{Pfn}}
\newcommand{\ran}{\operatorname{ran}}
\newcommand{\Dom}{\operatorname{Dom}}
\title{Local connectedness, range functors,\\
and essential incompleteness}
\author{Florian Lengyel}
\date{30 October 2026}
\begin{document}
\maketitle

\begin{abstract}
Using categories of partial computations, we prove that every consistent
computably enumerable extension of Peano arithmetic is incomplete.
The obstruction is a recursively inseparable pair: two disjoint
computably enumerable sets for which no recursive set contains the
first and avoids the second. We interpret programs in countable
models to obtain faithful functors from categories defined by provable
program equality. These functors preserve computation domains and
images. Each map selecting between two cases proved disjoint and
exhaustive by the theory is the interpretation of a program.
This verifies the hypotheses of a categorical
inseparability theorem. Finite computations then place a fixed
numerical inseparable pair inside the proof and refutation sets
of every consistent computably enumerable extension.
We also construct a classical propositional theory using computable
diagrams on a three-point space. Every consistent extension represented
by a computably enumerable subdiagram is incomplete.
\end{abstract}

\section{Introduction}

\subsection{Recursion categories and incompleteness}

Di Paola and Heller developed recursion categories as an algebraic
setting for recursion theory, with the aim of obtaining a categorical
formulation of incompleteness \cite{DPH}. Morphisms represent partial
computations, and restriction idempotents represent their domains.
A restriction idempotent is the partial identity that keeps precisely
the inputs on which a computation is defined.
Rosolini's $P$-categories provide a more general setting for these
constructions \cite{RosoliniThesis,Rosolini}. They describe partial
maps with a product operation. His representation theorem realizes such a
category inside a category of partial maps with a specified class of
monomorphisms as permitted domains \cite[Theorem 1.2]{Rosolini}.
Ranges express images, and coproducts provide the two branches used
in definition by cases.

One of the generalized incompleteness statements proposed by
Di Paola and Heller concerns recursive inseparability of theorems and
refutable sentences \cite[pp.~594--595]{DPH}.
For disjoint computably enumerable (c.e.) sets $A,B$, a \emph{recursive
separator} is a recursive set $C$ with $A\subseteq C$ and
$B\cap C=\varnothing$. The pair is \emph{recursively inseparable}
when no such $C$ exists. If the theorem and refutation sets of a
consistent computably enumerable theory contain such a pair,
the theory cannot be complete: searching simultaneously for a proof
and a refutation would decide a separator. When the same obstruction
holds in every consistent computably enumerable extension, the theory
is \emph{essentially incomplete}.

Bernardi developed an earlier categorical approach through precomplete
equivalence relations \cite{Bernardi}. In the numerical case,
precompleteness means that every partial recursive function agrees,
where defined and modulo the equivalence relation, with a total
recursive function. His categorical results establish inseparability
for nonempty disjoint unions of equivalence classes under the
hypotheses of his categorical theorem, and his examples include program equality
provable in Peano arithmetic \cite[pp.~125--128]{Bernardi}. This connects provable
program equality with inseparability. The construction below uses
domains, ranges, and local connectedness to obtain the obstruction,
then makes explicit its relation to proof and refutation.

\subsection{Local connectedness and comparison functors}

A computational domain is \emph{complemented} if it and a disjoint
computational domain cover their object. The categorical analogue of
an inseparable pair consists of disjoint domains that no complemented
domain separates. A domain is \emph{connected} if it cannot be
decomposed as the union of two disjoint nonzero domains.
A category is \emph{locally connected} when agreement of partial
arrows on every connected domain implies their equality.
Thus local connectedness supplies domains on which the behavior of
partial computations can be compared.

Heller used local connectedness in his existence theorem for recursion
categories \cite[Sections 8--9]{Heller}. In particular, it enters the
proof that computations built from iterated total operations are closed
under composition. The theorem constructs a recursion category from
the stated generation and iteration data; it does not require that
the ambient category already have a universal program.
Lengyel extended this existence argument by using faithful functors
that preserve domains, coproducts, zero maps, and iteration
\cite{Lengyel2004,Corrections}. The target must have enough connected
domains to detect equality. These are iteration functors; preservation
of products is not required in the corrected theorem.

The inseparability theorem of \cite{LOC} concerns a recursion category
that is already given. Its comparison is a \emph{range functor},
which preserves restriction domains, ranges, zero maps, and binary
coproducts. Under the target hypotheses stated in
Theorem~\ref{thm:local-lifting}, local connectedness constructs a
partial injection from two disjoint domains. The target also treats
domains as objects and admits partial inverses of partial injections, and its
inverse-image operations preserve binary unions of domains.
For complementary source domains, these properties yield a total
target map selecting one coproduct injection on the image of the
first domain and the other injection on the image of its complement.
The \emph{lifting hypothesis} requires this particular map to be the
image of a source arrow. Universality of computation in the source
then contradicts the existence of a separator.
Joins of source arrows with disjoint domains suffice for lifting
\cite[Proposition 4.8]{LOC}; the counterexample in
\cite[Section 4.5]{LOC} shows that the other hypotheses alone do not.

\subsection{Syntactic categories and the present applications}

Syntactic recursion categories explain why local connectedness is
imposed on the target of a functor rather than on its source.
For a consistent extension $T$ of Peano arithmetic ($\PA$), the category
$S_T$ has programs as arrows, identified when $T$ proves one
universally quantified sentence asserting equality of their partial
graphs \cite{Montagna,DPM}. Montagna's category $S$ is the case
$T=\PA$. Di Paola and Montagna prove that $S_T$ is not locally
connected for any consistent computably enumerable $T$
\cite[Corollary 4.4]{DPM}. Their examples therefore motivate
comparison with a category in which connected domains do determine
arrows, as in the existence and inseparability results just described.

The primed systems use a different equivalence. Montagna's $S'$ is
the case $T=\PA$ of the system $S'_T$ studied by Di Paola and
Montagna: two programs are identified if $T$ proves equality of their
partial computations separately at each standard input.
The claimed local-connectedness result for these systems
\cite[Theorem 4.5, pp.~658--659]{DPM} includes an arithmetic
soundness hypothesis. However, \emph{Pointwise provable equality and
the failure of composition} \cite[Corollary 2]{Composition} proves
that their proposed composition is not independent of representatives
for any consistent computably enumerable extension $T$ of $\PA$.
Thus $S'$ and $S'_T$ do not define categories under that composition
and cannot supply locally connected examples here.
The single uniformly quantified proof used in the unprimed $S_T$
makes composition well-defined; its verification is included in
Section~\ref{sec:syntactic}.

We construct such a comparison by interpreting programs in all
countable models of $T$. The interpretations form a faithful range
functor to a product of categories of partial functions, and that
product is locally connected. Given complementary source domains,
one program searches their defining predicates and returns the
appropriately tagged input. The theory proves that the program is
total; its interpretation is the required selecting map in every
model. Restricting the model family makes the construction compatible
with theory extension. This verifies the full lifting hypothesis
and gives an inseparable pair in each $S_T$.

The arithmetic conclusion requires an effective observation of this
construction. The same local-connectedness theorem supplies a fixed
inseparable pair in the classical category of partial recursive
functions. Finite computations give PA-proofs of the corresponding
numeral instances. Consequently the proof and refutation sets of
every consistent computably enumerable extension of $\PA$ contain
that numerical pair. Completeness would produce a recursive
separator, proving essential incompleteness without a soundness
assumption. The faithful model interpretation and the effective
numeral observations have distinct roles: the former represents the
syntactic category, while the latter connects computation with
arithmetical proof.

A second application uses the computational cosheaves of
\cite{ThreePoint}. For the three-point space considered there,
these are computable diagrams consisting of two arrows with a
common source. Their partial-map category contains a nonnumerical
recursion category and a classical propositional presentation.
Taking the component at either target of the diagram gives a
functor to numerical partial computations. We prove that it preserves
the chosen universal computation and use the proofs of atomic
axioms to obtain an inseparable pair within each extension's
proof and refutation sets. The extensions here are the rule-closed
computational domains of the constructed syntax object. Every
consistent such extension is incomplete.

Section~\ref{sec:preliminaries} fixes the categorical definitions,
and Section~\ref{sec:syntactic} defines $S_T$.
Sections~\ref{sec:models} and \ref{sec:family} construct its
model interpretation and the comparison under theory extension.
Section~\ref{sec:lifting} proves the needed local-connectedness
results and verifies the lifting hypothesis.
Section~\ref{sec:arithmetic} obtains arithmetic essential
incompleteness, and Section~\ref{sec:cosheaves} develops the
three-point application.

\section{Categorical preliminaries}\label{sec:preliminaries}

\subsection{Partial maps and domains}

Composition is written $gf=g\circ f$, so $f$ is applied first.
A \emph{restriction category} assigns to every arrow
$f:X\rightharpoonup Y$ an arrow
$\bar f:X\rightharpoonup X$ satisfying
\[
 f\bar f=f,\qquad
 \bar f\bar g=\bar g\bar f,\qquad
 \overline{g\bar f}=\bar g\bar f,\qquad
 \bar h f=f\overline{hf}.
\]
In the middle two equations $g:X\rightharpoonup Z$; in the last,
$h:Y\rightharpoonup Z$ \cite{CockettLack}.
The arrows $\bar f$ are idempotent. They are called
\emph{restriction idempotents}, or \emph{domains}.
The arrow $f$ is \emph{total} when $\bar f=1_X$.
The total arrows form a subcategory, denoted by
$\operatorname{Tot}(\mathcal C)$ for a restriction category $\mathcal C$.

Domains on $X$ are ordered by $a\le b$ when $ba=a$.
Their meet is composition, their greatest element is $1_X$, and
$a\vee b$ denotes their least upper bound when it exists.
For parallel partial arrows, the restriction order is
$f\le g$ when $f=g\bar f$.
A \emph{restriction zero} $0_{X,Y}:X\rightharpoonup Y$ is
absorbing under composition and satisfies
$\overline{0_{X,Y}}=0_{X,X}$. We abbreviate $0_{X,X}$ to $0_X$.
Domains $a,b$ on $X$ are disjoint when $ab=0_X$.
They are complementary when also $a\vee b=1_X$.
A disjoint pair $a_0,a_1$ is \emph{inseparable} if no complementary
pair $c,c'$ satisfies $a_0\le c$ and $a_1\le c'$.

For $f:X\rightharpoonup Y$ and a domain $a$ on $Y$, its
\emph{inverse image} is the domain $f^*a=\overline{af}$ on $X$.
A domain $r$ on $Y$ \emph{receives} $f$ when $rf=f$.
The \emph{range} $\ran(f)$ is the least domain receiving $f$;
we require the agreement law
\[
 gf=hf\quad\Longrightarrow\quad
 g\ran(f)=h\ran(f)
 \qquad(g,h:Y\rightharpoonup Z).
\]
For ordinary partial functions, domains and ranges are the partial
identities on their domains of definition and images.

A domain $e$ on $X$ \emph{splits} if
$e=mr$ for a total $m:A\to X$ and $r:X\rightharpoonup A$
with $rm=1_A$. The category is \emph{split} when every domain splits.
Its \emph{restriction splitting} has objects $(X,e)$, where $e$ is a
domain on $X$, and arrows
$f:(X,e)\rightharpoonup(Y,d)$ given by original arrows
$f:X\rightharpoonup Y$ satisfying $f=dfe$.
The identity of $(X,e)$ is $e$, and $f$ is total there when
$\bar f=e$. The object $(X,e)$ thus represents the domain as an
object in the corresponding total category.

\subsection{Products, coproducts, and local connectedness}

A \emph{binary restriction product} $X\times Y$ has total
projections $\pi_1:X\times Y\to X$ and $\pi_2:X\times Y\to Y$.
For arrows $f:Z\rightharpoonup X$ and $g:Z\rightharpoonup Y$,
there is a unique $\langle f,g\rangle:Z\rightharpoonup X\times Y$
with
\[
 \pi_1\langle f,g\rangle=f\bar g,\qquad
 \pi_2\langle f,g\rangle=g\bar f,\qquad
 \overline{\langle f,g\rangle}=\bar f\bar g.
\]
Products of arrows are obtained by this pairing and the projections.
The total categorical diagonal is
$\Delta_X=\langle1_X,1_X\rangle:X\to X\times X$.
This is the restriction-category description of the product operation
in a $P$-category \cite{Rosolini,CockettLack}.

Binary \emph{restriction coproducts} are categorical coproducts with
total injections $i_0:X\to X+Y$ and $i_1:Y\to X+Y$.
For arrows $f:X\rightharpoonup Z$ and $g:Y\rightharpoonup Z$,
write $[f,g]:X+Y\rightharpoonup Z$ for their copairing.
The \emph{codiagonal} is
$\nabla_X=[1_X,1_X]:X+X\to X$.
In a category with ranges and restriction coproducts,
$a\vee b=\ran[a,b]$ for domains on the same object.

A \emph{range functor} $F:\mathcal C\to\mathcal D$ is a covariant
functor preserving restrictions, ranges, restriction zeros, and binary
restriction coproducts, including their injections and copairings
under the coproduct comparison isomorphisms. It preserves binary
domain joins by the preceding formula, and hence complementary pairs.

A domain $a$ is \emph{connected} if $a=b\vee c$ and $bc=0$
imply $b=0$ or $c=0$. The category is \emph{locally connected}
when its connected domains are jointly epimorphic: for
$f,g:X\rightharpoonup Y$, the equations $fa=ga$ for every connected
domain $a$ on $X$ imply $f=g$.
A \emph{partial monomorphism} $p:X\rightharpoonup Y$ satisfies
$pu=pv\Rightarrow\bar p u=\bar p v$ for all
$u,v:W\rightharpoonup X$.
A \emph{section} of $p$ is an arrow $s:Y\rightharpoonup X$
with $ps=\bar s$ and $psp=p$.
\emph{Weak choice} means that every partial monomorphism has such
a section. For partial functions, this is the existence of the
partial inverse of an injection on its domain.

\subsection{Recursion categories}

A restriction category with binary restriction products and zeros is
\emph{prodominical} if $f\times0=0$ for every arrow $f$ and every
zero arrow of the required type. It is \emph{isotypical} if any
two of its objects are isomorphic.

A \emph{Turing morphism} on $U$ is an arrow
$\tau:U\times U\rightharpoonup U$ such that, for every object $X$
and every $f:X\times U\rightharpoonup U$, there is a total
parameter map $r:X\to U$ satisfying
\[
 f=\tau(r\times1_U).
\]
A \emph{recursion category} here is a prodominical isotypical
category with binary restriction products and a Turing morphism,
as in \cite[Section 9]{Heller}.
Ranges and coproducts are stated as additional hypotheses when used.

\section{The syntactic recursion category}\label{sec:syntactic}

Let $T$ be a consistent classical theory in the language of arithmetic
containing $\PA$. Write $\N$ for the natural numbers and $\omega$
for their formal numerical type. The other formal numerical types are
generated from $\omega$ by binary products and tagged coproducts.
Fix recursive numerical codings
whose bijection and structural equations are provable in $\PA$.
For a typed program index $e:X\rightharpoonup Y$, let $G_e(x,y)$ be its
$\Sigma_1$ output formula. It existentially quantifies a computation trace
over a primitive-recursive trace test. Fix a deterministic programming
system whose graph identities, including functionality, composition,
and parameterization, are provable in $\PA$.

The category $S_T$ has these formal types as objects and typed indices
modulo the equivalence relation
\begin{equation}\label{eq:uniform}
 e\equiv_T f
 \quad\Longleftrightarrow\quad
 T\vdash \forall x\forall y\,
       \bigl(G_e(x,y)\leftrightarrow G_f(x,y)\bigr).
\end{equation}
The quantifiers belong to the single sentence proved in $T$.
This is the unprimed syntactic construction of Montagna and
Di Paola--Montagna \cite{Montagna,DPM}.

For a program operation $\operatorname{comp}$ implementing composition,
$\PA$ proves
\begin{equation}\label{eq:composition}
 G_{\operatorname{comp}(f,e)}(x,z)
 \leftrightarrow
 \exists y\,\bigl(G_e(x,y)\land G_f(y,z)\bigr).
\end{equation}
Substitution of uniformly equivalent graphs under this existential proves
that composition descends to \eqref{eq:uniform}. The identity and
associativity equations follow from the same graph identities.

Every $\Sigma_1$ formula $C(x)$ represents a partial identity: search for
a witness to $C(x)$ and then return $x$. Write $[C]_T$ for this domain.
Then
\begin{equation}\label{eq:domain-order}
 [C]_T\le[D]_T
 \quad\Longleftrightarrow\quad
 T\vdash\forall x\,(C(x)\longrightarrow D(x)).
\end{equation}
Conjunction and disjunction give domain meets and joins. The domain and
range of $[e]_T$ are represented respectively by
\[
 \exists y\,G_e(x,y),\qquad
 R_e(y):=\exists x\,G_e(x,y).
\]
For the range property, a domain $[C]_T$ receives $[e]_T$ precisely when
$T$ proves $G_e(x,y)\to C(y)$ universally, equivalently when
$T$ proves $R_e(y)\to C(y)$ universally. For agreement, a witness to
$R_e(y)$ together with equality of two composite graphs gives equality
of the two output graphs at $y$. Existential elimination in $T$ proves
the required agreement equation. These arguments also show that the
chosen ranges are unchanged under theory extension.

The ordinary product and tagged-coproduct programs supply the restriction
products and coproducts. The universal program supplies a Turing morphism
\[
 \tau:\omega\times\omega\rightharpoonup\omega:
 \quad f:X\times\omega\rightharpoonup\omega
 \ \Longrightarrow\ 
 f=\tau(r\times1_\omega)
\]
for some total parameter arrow $r:X\to\omega$. Formalized
parameterization makes this arrow $\PA$-provably total. The chosen coding
isomorphisms make $S_T$ isotypical, and product with a zero arrow is zero.
Thus $S_T$ is a prodominical recursion category. The preceding
constructions do not require that the axioms of $T$ be computably enumerable.

\section{Interpretation in models}\label{sec:models}

The category $\Pfn$ has sets as objects and partial functions as arrows,
with ordinary composition. Restrictions and ranges are partial identities
on domains of definition and images. Restriction products are Cartesian products;
coproducts are tagged disjoint unions.

Fix a model $M\models T$. Interpret numerical types recursively by
\[
 X_M=|M|\quad(X=\omega),\qquad
 (X\times Y)_M=X_M\times Y_M,\qquad
 (X+Y)_M=X_M\amalg Y_M.
\]
The fixed arithmetic codings identify these sets with the corresponding
codes in $M$. A standard program index in a graph formula is interpreted
as its numeral in $M$. Define
\begin{equation}\label{eq:model-functor}
 F_M(X)=X_M,\qquad
 F_M([e]_T)(x)=y
 \quad\Longleftrightarrow\quad M\models G_e(x,y).
\end{equation}
All quantified traces and intermediate values in this formula range over
the whole model, including its nonstandard elements.

\begin{proposition}\label{prop:model}
For every $M\models T$, \eqref{eq:model-functor} defines a covariant
functor $F_M:S_T\to\Pfn$ preserving restrictions, zeros, binary
restriction products and coproducts, ranges, and finite joins of domains.
\end{proposition}
\begin{proof}
Functionality is provable in $\PA$, so the interpreted graph is a
partial function. Equation \eqref{eq:uniform} makes the assignment
independent of the representative. Equation \eqref{eq:composition},
interpreted in $M$, is exactly relational composition of the two partial
functions; their intermediate values range over $Y_M$. The identity
graph gives the identity function.

The domain predicate $\exists y\,G_e(x,y)$ gives the interpreted
function's domain. Similarly,
\[
 \{y\in Y_M:M\models R_e(y)\}
 =\operatorname{im}(F_M[e]_T),
\]
so ranges are preserved. False predicates, conjunction, and disjunction
give zeros, intersections, and unions. The PA-provable product and
tagged-coproduct graph equations give the corresponding operations on
partial functions.
\end{proof}

\section{A faithful family and theory extension}\label{sec:family}

Let $\mathcal M_T$ be the set of all structures for the language of
arithmetic with underlying set $\N$ that satisfy $T$. Their arithmetic
operations need not be the standard ones, and the interpretation of a
numeral $\bar n$ need not be the carrier element $n$. Completeness and
the downward L\"owenheim--Skolem theorem imply
$\mathcal M_T\ne\varnothing$: every model of $T$ is infinite, and a
countably infinite model can be transported to the carrier $\N$.

Define the product category
\[
 \mathcal D_T=\prod_{M\in\mathcal M_T}\Pfn.
\]
An object is a family of sets $(A_M)_M$; an arrow is a family of partial
functions with the indicated sources and targets. Composition and the
restriction, range, restriction-product, and coproduct operations are coordinatewise.
The functors of Proposition~\ref{prop:model} combine to give
\[
 F_T:S_T\longrightarrow\mathcal D_T,\qquad
 F_T(X)=(X_M)_M,\quad F_T([e]_T)=(F_M[e]_T)_M.
\]

\begin{theorem}\label{thm:faithful}
The functor $F_T$ is faithful and preserves restrictions, zeros, binary
restriction products and coproducts, ranges, and finite domain joins. The category
$\mathcal D_T$ is split and locally connected, satisfies weak choice,
and has inverse images preserving binary domain joins. Moreover,
$F_T(1_\omega)\ne0$.
\end{theorem}
\begin{proof}
Preservation follows coordinatewise from Proposition~\ref{prop:model}.
If $[e]_T\ne[f]_T$, let $E$ denote the uniform graph-equality sentence
in \eqref{eq:uniform}. The theory $T\cup\{\neg E\}$ is consistent.
Completeness and downward L\"owenheim--Skolem give a model on carrier
$\N$ in which the two graphs differ. This is a coordinate of $F_T$,
which therefore distinguishes the arrows.

A domain in $\mathcal D_T$ is a family of subsets. Its splitting uses
those subsets as the coordinate sets. For an object $A=(A_M)_M$, a
domain supported at one coordinate $M$, with a singleton subset
$\{x\}\subseteq A_M$ there, is connected. These domains detect the
definedness and value of every coordinate of a partial arrow. Hence
$\mathcal D_T$ is locally connected.

A partial monomorphism has injective coordinate functions: test its
defining property on singleton objects supported at one coordinate.
The partial inverses of these injections form a section, proving weak
choice. Inverse images commute with unions coordinatewise. Finally,
$\mathcal M_T$ is nonempty and every interpreted numerical object is
nonempty, so its identity differs from zero.
\end{proof}

For consistent computably enumerable $T\supseteq\PA$, the source $S_T$
is not locally connected \cite[Corollary 4.4, pp.~657--658]{DPM}.
The connected domains used in the preceding proof belong to
$\mathcal D_T$ and need not be images of source domains. Faithfulness
reflects equality of arrows; it does not require these additional target
domains to have source preimages.

For an inclusion $T\subseteq T'$ of consistent theories, define
\[
 Q_{T,T'}:S_T\to S_{T'},\qquad
 Q_{T,T'}X=X,\quad Q_{T,T'}[e]_T=[e]_{T'}.
\]
Provability is monotone, so this is a well-defined covariant functor.
It preserves the chosen computational operations and Turing morphism.
Since $\mathcal M_{T'}\subseteq\mathcal M_T$, restriction of families
defines a covariant projection
$\pi_{T,T'}:\mathcal D_T\to\mathcal D_{T'}$. Fixing type interpretations
once for each arithmetic structure gives a strictly commuting square:
\[
\begin{tikzcd}[column sep=large]
 S_T \arrow[r,"Q_{T,T'}"] \arrow[d,"F_T"'] &
 S_{T'} \arrow[d,"F_{T'}"]\\
 \mathcal D_T \arrow[r,"\pi_{T,T'}"'] & \mathcal D_{T'}.
\end{tikzcd}
\]
Both paths send a program to the family of its interpretations in the
models of $T'$. The identity and composition laws for the projections
and quotient functors follow from restriction of indices and the
unchanged program representatives.

\section{Switching and the local-connectedness theorem}\label{sec:lifting}

Theorem~\ref{thm:faithful} supplies
such a functor into a split locally connected category satisfying weak
choice and inverse-image preservation of domain joins. These are the
target hypotheses of the local-connectedness inseparability theorem
\cite[Theorem 4.7]{LOC}. Its remaining requirement concerns lifting
specific sections associated with complemented source domains.

\subsection{The local-connectedness construction}

The following argument supplies the needed sections. It is the
construction of \cite[Section 4.3]{LOC}, with its restriction-category
hypotheses stated explicitly. For a coproduct $Y+Z$, let
$\varepsilon_0=1_Y+0_Z$ and $\varepsilon_1=0_Y+1_Z$ be its summand
domains. Their meet is zero and their join is the identity. These facts
follow by precomposition with the two injections and the agreement
property of ranges.

\begin{lemma}\label{lem:connected-alternative}
Let $\mathcal A$ be a restriction category with ranges, restriction zeros,
and binary restriction coproducts. Suppose inverse images preserve binary
domain joins. If $a$ is a connected domain on $X$ and
$f:X\to Y+Z$ is total, then
$\varepsilon_0fa=fa$ or $\varepsilon_1fa=fa$.
\end{lemma}
\begin{proof}
Totality and the inverse-image laws give the disjoint decomposition
\[
 a=af^*\varepsilon_0\vee af^*\varepsilon_1.
\]
Connectedness makes one term zero and the other equal to $a$. In the
latter case $a\le f^*\varepsilon_i$, and the restriction identity
$\varepsilon_i f=f\overline{\varepsilon_i f}$ gives
$\varepsilon_i fa=fa$.
\end{proof}

\begin{proposition}\label{prop:local-switch}
Let $\mathcal A$ be a locally connected split restriction category with
ranges, restriction zeros, and binary restriction coproducts. Suppose
weak choice holds and inverse images preserve binary domain joins.
For disjoint domains $a,b$ on $X$, the arrow
$p=\nabla_X(a+b):X+X\rightharpoonup X$ is a partial monomorphism
with a unique section $s:X\rightharpoonup X+X$, and
\[
 ps=\bar s=a\vee b.
\]
In particular, for complementary $c,c'$ on $X$, the arrow
$\nabla_X(c'+c)$ has a total section $s:X\to X+X$ satisfying
$sc=i_1c$ and $sc'=i_0c'$.
\end{proposition}
\begin{proof}
First let $\theta,\theta':Y\to X+X$ be total and suppose
$p\theta=p\theta'$. The identities
$ap=\nabla_X(a+0)$ and $bp=\nabla_X(0+b)$ give
\[
 \nabla_X(a+0)\theta=\nabla_X(a+0)\theta',
 \qquad
 \nabla_X(0+b)\theta=\nabla_X(0+b)\theta'.
\]
For a connected domain $\gamma$ on $Y$,
Lemma~\ref{lem:connected-alternative} places each of
$\theta\gamma,\theta'\gamma$ in one summand.
If both lie in the first summand, the first equality, followed by $i_0$,
gives $(a+b)\theta\gamma=(a+b)\theta'\gamma$; the second-summand
case uses the second equality and $i_1$.
If they lie in opposite summands, the two equalities force both these
restricted arrows to be zero. Local connectedness therefore gives
$(a+b)\theta=(a+b)\theta'$. Since $\bar p=a+b$, this proves the
partial-monomorphism condition for total test arrows.

For arbitrary $u,v:W\rightharpoonup X+X$ with $pu=pv$, split their
common composite domain as
$\overline{pu}=\overline{pv}=me$, where $m:Z\to W$ is total and
$em=1_Z$. The arrows $pum,pvm$, hence also $um,vm$, are total.
The preceding case gives $\bar p um=\bar p vm$. Composing with $e$
and using $\bar p u=\bar p u\,\overline{pu}$ and its analogue for $v$
gives $\bar p u=\bar p v$. Thus $p$ is a partial monomorphism.

Weak choice supplies a section $s$. The equations $ps=\bar s$ and
$psp=p$ imply $\bar s=\ran(p)$: this domain receives $p$, and every
domain receiving $p$ also receives $ps=\bar s$. Moreover
$\bar p s=s\overline{ps}=s$. Two sections have the same composite
with $p$, so partial monicity makes them equal.
The least-domain property gives $\ran(p)=a\vee b$.

For $a=c'$ and $b=c$, the section is total because $c'\vee c=1_X$.
The equality $p(sc)=c=p(i_1c)$ and partial monicity give $sc=i_1c$;
the equation on $c'$ is analogous.
\end{proof}

\begin{theorem}[Range-functor lifting]\label{thm:local-lifting}
Let $\mathscr R$ be a prodominical recursion category with ranges and
binary restriction coproducts. Let $\mathcal A$ be a locally connected
split restriction category with ranges, restriction zeros, and binary
restriction coproducts, satisfying weak choice, with inverse images
preserving binary domain joins. Let
$F:\mathscr R\to\mathcal A$ be a range functor.
Choose a Turing morphism $\tau:U\times U\rightharpoonup U$ and
an isomorphism $b:U\to U+U$, with $F(1_U)\ne0$.
For each complementary pair $c,c'$ on $U$, let $s_c$ be the section
of $\nabla_{FU}(Fc'+Fc)$. Suppose there is an arrow
$\ell_c:U\rightharpoonup U+U$ with $F\ell_c=s_c$, using the coproduct
comparison. Then
\[
 a_0=(b\tau\Delta_U)^*(1_U+0_U),\qquad
 a_1=(b\tau\Delta_U)^*(0_U+1_U)
\]
are inseparable in $\mathscr R$.
\end{theorem}
\begin{proof}
This is \cite[Theorem 4.7]{LOC}; the following proof uses the
inverse-image laws directly after constructing the section.
The domains $a_0,a_1$ are disjoint. Suppose $c,c'$ separate them.
Put $u=b^{-1}\ell_c$, and choose a total parameter $g:U\to U$ with
$\tau(g\times1_U)=u\pi_2$. Precomposition with $\langle1_U,g\rangle$
gives $\tau\Delta_Ug=ug$.

In the target write $E=Fg$, $H=F(b\tau\Delta_U)$, $\Gamma=Fc$,
and $\Gamma'=Fc'$. The preceding equation gives $HE=s_cE$.
Since $s_c\Gamma=i_1\Gamma$ and $s_c\Gamma'=i_0\Gamma'$,
its inverse images of the summand domains are $\Gamma'$ and $\Gamma$,
respectively. Preservation of restrictions and composition now gives
\[
 E^*F(a_0)=E^*\Gamma',\qquad E^*F(a_1)=E^*\Gamma.
\]
The separator inequalities imply that $E^*\Gamma$ and $E^*\Gamma'$
are equal. They are complementary because $E$ is total and target
inverse images preserve joins. Hence $1_{FU}=0$, contradicting the
hypothesis. The proof uses totality of $F\ell_c=s_c$, not totality of
its source lift $\ell_c$.
\end{proof}

\subsection{Lifting in the syntactic category}

\begin{lemma}\label{lem:switch}
Let $c,c'$ be complementary domains on $X$ in $S_T$. There is a total
arrow $s:X\to X+X$ with
\[
 sc=i_1c,\qquad sc'=i_0c'.
\]
Its image under $F_T$ is the unique section of
$\nabla_{F_TX}(F_Tc'+F_Tc)$, using the coproduct comparison to identify
$F_T(X+X)$ with $F_TX+F_TX$.
\end{lemma}
\begin{proof}
Choose $\Sigma_1$ representatives $C(x)=\exists u\,R(x,u)$ and
$C'(x)=\exists v\,R'(x,v)$ with primitive-recursive witness tests.
The theory $T$ proves universally their disjointness and coverage.
Search increasing finite bounds for a witness to either predicate,
returning the tagged input $i_1x$ on finding a $C$-witness and $i_0x$
on finding a $C'$-witness. Fix a priority for simultaneous discoveries.
The finite-search verification in $\PA$ proves that the program halts
exactly when $C(x)\lor C'(x)$. Under the disjointness assumption it also
proves the exact graph equation
\[
 G_s(x,z)\leftrightarrow
 (C(x)\land z=i_1x)\lor(C'(x)\land z=i_0x).
\]
Thus $T$ proves this equation and universal halting. The equations for
the restrictions follow, giving a total arrow of the uniform quotient.

In each $M\models T$, the interpretations of $C,C'$ partition $X_M$.
The interpreted program sends each input to its specified tagged copy.
It is exactly the inverse of the partial injection
$\nabla_{X_M}(c'_M+c_M)$. Its graph is therefore the unique required
section. The same program works at every coordinate, proving the lifting
statement for the entire family.
\end{proof}

\begin{corollary}\label{cor:pair}
For every consistent $T\supseteq\PA$, the functor $F_T$ satisfies the
lifting hypothesis of the local-connectedness inseparability theorem.
Fix $U=\omega$, a Turing morphism $\tau:U\times U\rightharpoonup U$,
and a PA-provable coding isomorphism $b:U\to U+U$. Put
\[
 d=\tau\Delta_U,\qquad
 a_0=(bd)^*(1_U+0_U),\qquad a_1=(bd)^*(0_U+1_U).
\]
Then $a_0,a_1$ are inseparable in $S_T$. Their images under
$Q_{T,T'}$ are the identically defined inseparable pair in $S_{T'}$.
\end{corollary}
\begin{proof}
Theorem~\ref{thm:faithful} and Lemma~\ref{lem:switch} verify all the
hypotheses of Theorem~\ref{thm:local-lifting}. The quotient functor keeps
every program in the definition fixed, and the same argument applies
to consistent $T'$.
\end{proof}

The local-connectedness theorem uses Turing universality in its source.
The target $\mathcal D_T$ has all families of partial functions as arrows;
$F_T(\tau)$ is not universal for these arrows: at each model coordinate,
countably many program parameters cannot enumerate all partial functions
on its countable numerical carrier. The target's domain lattices are
Boolean, so the image of every disjoint source pair has a separator
there. Lemma~\ref{lem:switch} lifts the sections associated with source
complemented domains. An arbitrary target separator need not come from
such a domain. This distinction explains how a faithful range functor
can have an inseparable source pair and a separable image pair.

\section{Effective observations and essential incompleteness}\label{sec:arithmetic}

Let $\mathrm{Rec}$ be the classical recursion category on the formal
numerical types, with partial recursive functions as arrows and extensional
equality. Its domains are computably enumerable subsets, and its
complemented domains are recursive subsets.
Assume now that $T$ has computably enumerable axioms. For a domain
$c=[C]_T$ on $\omega$, define the subset of the standard natural numbers
\begin{equation}\label{eq:observation}
 E_T(c)=\{n\in\N:T\vdash C(\bar n)\}.
\end{equation}
It is computably enumerable by proof enumeration and is independent of
the representative. The operation is monotone, preserves finite meets,
and preserves top and bottom. Let $\kappa_n:\omega\to\omega$ be the
total constant program with value $n$. Then
\begin{equation}\label{eq:faithful-observation}
 n\in E_T(c)
 \quad\Longleftrightarrow\quad c\kappa_n=\kappa_n
 \quad\Longleftrightarrow\quad
 F_T(c)F_T(\kappa_n)=F_T(\kappa_n).
\end{equation}
The first equivalence follows from the partial-identity graph equation;
the second is faithfulness. Thus the standard observations can also be
expressed by equations in the faithful representation.

Take $b(2m)=i_0m$ and $b(2m+1)=i_1m$ in
Corollary~\ref{cor:pair}, using an acceptable universal program.
Let $\alpha(x),\beta(x)$ be the resulting $\Sigma_1$ domain formulas.
Their standard truth sets
\[
 A=\{n:\N\models\alpha(\bar n)\},\qquad
 B=\{n:\N\models\beta(\bar n)\}
\]
are the classical even-output and odd-output sets for self-application.
Their recursive inseparability is another instance of
Theorem~\ref{thm:local-lifting}. To see this, let
$J:\mathrm{Rec}\to\Pfn$ send a formal numerical type to its underlying
set and a partial recursive arrow to the same partial function.
It is a covariant range functor. The category $\Pfn$ satisfies
the target hypotheses: singleton domains detect partial functions,
domains split as subsets, partial injections have partial inverses,
and inverse images preserve unions. Complementary domains in
$\mathrm{Rec}$ are recursive sets. Their deciding procedure computes
the required tagged-input map, so its underlying partial function is
the target section. This supplies the lifting hypothesis.
The theorem, with the chosen universal program and parity
isomorphism $b$, gives precisely the pair $A,B$.
Each finite computation gives a PA-proof of its numeral instance, and $\PA$ proves
$\forall x\,\neg(\alpha(x)\land\beta(x))$.

Define
\[
 P_T=\{n:T\vdash\alpha(\bar n)\},\qquad
 N_T=\{n:T\vdash\neg\alpha(\bar n)\}.
\]
These are computably enumerable disjoint sets with $A\subseteq P_T$
and $B\subseteq N_T$. They enlarge under theory extension.
Let $\operatorname{Pair}(\mathrm{Rec})$ have objects $(X,A,B)$ consisting
of a numerical type and disjoint computably enumerable domains.
An arrow $(X,A,B)\to(Y,P,N)$ is a total recursive function $f:X\to Y$
with $f(A)\subseteq P$ and $f(B)\subseteq N$. Identities and composition
are inherited from $\mathrm{Rec}$. The assignment
$T\mapsto(\omega,P_T,N_T)$, with the identity function for theory
inclusions, is a covariant functor from the poset of consistent
computably enumerable extensions of $\PA$ to
$\operatorname{Pair}(\mathrm{Rec})$.
The inclusions of $A,B$ give identity-arrow reductions to every such pair.
A recursive separator pulls back along a total recursive function;
consequently $P_T,N_T$ are recursively inseparable.

\begin{corollary}[Essential incompleteness of arithmetic]\label{cor:arithmetic}
Every consistent computably enumerable extension of $\PA$ is incomplete.
\end{corollary}
\begin{proof}
If $T$ were complete, parallel enumeration of proofs of
$\alpha(\bar n)$ and $\neg\alpha(\bar n)$ would halt on every standard
input. Consistency makes the answers exclusive. The resulting recursive
set $P_T$ would separate $A$ and $B$, a contradiction.
\end{proof}

The effective step uses standard numeral instances and the presentation
of $T$ by computably enumerable axioms. The semantic functor exists also
for $T=\operatorname{Th}(\N)$, when $T$ is complete and $S_T$ is the
classical category of partial recursive functions on the formal numerical
types. Thus internal inseparability alone does not imply logical
incompleteness. Bernardi's precomplete-equivalence method gives another
established route through PA-provable program equivalence
\cite[pp.~125--128]{Bernardi}.

\section{Computational cosheaves and effective observation}\label{sec:cosheaves}

For a topological space $X$, let $\mathcal O(X)$ be the category
of open sets and inclusions. A \emph{set-valued cosheaf} is a
covariant functor $C:\mathcal O(X)\to\mathbf{Set}$ for which
every open cover $(O_\lambda)_\lambda$ of $O$ gives a coequalizer
\[
 \coprod_{\lambda,\mu} C(O_\lambda\cap O_\mu)
 \rightrightarrows \coprod_\lambda C(O_\lambda)
 \longrightarrow C(O).
\]
The two parallel maps are induced by the inclusions into the two
members of the cover; the empty cover requires
$C(\varnothing)=\varnothing$.
Now let $X=\{a,b,c\}$ have opens
$\varnothing,\{c\},\{a,c\},\{b,c\},X$.
Its set-valued cosheaves are equivalent to covariant diagrams of shape
$J=(1\leftarrow0\to2)$: the three components are the values at
$\{c\},\{a,c\},\{b,c\}$, and the value at $X$ is their pushout
\cite[Section 2]{ThreePoint}.
Let $\mathcal C_{\mathrm{rec}}$ have computably enumerable subsets
of finite powers of $\N$ as objects and partial recursive functions
total on their sources as arrows, with extensional equality.
The computational diagrams form the total category
$\mathcal E=\mathcal C_{\mathrm{rec}}^J$.
An object $Y$ consists of arrows $r_i^Y:Y_0\to Y_i$, $i=1,2$,
and a morphism is a commuting triple of total recursive component maps.
Define the partial category
\[
 \mathcal A=\operatorname{Par}(\mathcal E,\operatorname{Mono}).
\]
Its arrows $Y\rightharpoonup Z$ are isomorphism classes of spans
$Y\xleftarrow{m}D\xrightarrow{f}Z$ with $m$ monic; composition
uses pullback. Monomorphisms are componentwise injections.
Thus a partial arrow is equivalently a computably enumerable
subdiagram $D\subseteq Y$ with recursive value maps $D_j\to Z_j$
commuting with the transitions. Its restriction is the partial
identity on $D$, and its range is the subdiagram of component images.

The diagram used for universal computation in
\cite[Sections 4--5]{ThreePoint} is
\begin{equation}\label{eq:universal-fork}
 V_0=\N^4,\qquad V_1=V_2=\N^2,\qquad
 r_1(k,z,a,b)=(k,a),\quad r_2(k,z,a,b)=(k,b).
\end{equation}
Every diagram $Y$ admits a total map $m:Y\to V$ with a partial
inverse $r:V\rightharpoonup Y$ such that $rm=1_Y$ and
$mr$ is a restriction idempotent. In this sense $V$ is a
\emph{universal object}.
Write $V^\#$ for the diagram of nonempty finite words in its
components, with transitions acting letter by letter.
The chosen \emph{frames} are computable isomorphisms
$V\cong V\times V$, $V\cong V+V$, and $V\cong V^\#$.
Let $\mathscr R_X$ denote the full subcategory of $\mathcal A$
on the formal objects generated by $V$ using those operations.
It is an isotypical recursion category, and its restriction
splitting is equivalent to $\mathcal A$.

\subsection{Local connectedness and lifting}

The subdiagram generated by a central element and its two endpoint
images is finite and connected. Singleton subdiagrams at either endpoint
are connected as well. All are computational domains. They detect every
component domain and value of a partial arrow, so $\mathcal A$ is locally
connected. Its splitness, weak choice, ranges, coproducts, and
componentwise inverse-image laws are established in
\cite[Proposition 2.2]{ThreePoint}.
The full inclusion $\mathscr R_X\hookrightarrow\mathcal A$ is a
range functor.
The section associated with a complemented domain has source and target
in $\mathscr R_X$; fullness therefore supplies its lift. This gives the
inseparability theorem in a nonnumerical prodominical recursion category.
Indeed, the terminal object of $\mathcal E$ has five subobjects,
whereas a numerical singleton has two.
A restriction equivalence with the classical recursion category
would extend to the restriction splittings and then restrict to an
equivalence $\mathcal E\simeq\mathcal C_{\mathrm{rec}}$, contradicting
this invariant. Moreover, $V$ has no decidable equality
\cite[Proposition 6.1]{ThreePoint}.

\subsection{Endpoint functors}

Let
$\mathcal K=\operatorname{Par}(\mathcal C_{\mathrm{rec}},\operatorname{Mono})$.
This is the restriction splitting of the classical recursion category:
objects are c.e.\ numerical carriers and arrows are partial recursive
functions with c.e.\ domains contained in their sources.
For $j\in\{0,1,2\}$, component evaluation defines
\begin{equation}\label{eq:endpoint-functor}
 E_j:\mathcal A\longrightarrow\mathcal K,\qquad
 E_j(Y)=Y_j,\qquad
 E_j(Y\xleftarrow{m}D\xrightarrow{f}Z)
   =(Y_j\xleftarrow{m_j}D_j\xrightarrow{f_j}Z_j).
\end{equation}

\begin{proposition}\label{prop:endpoint}
The covariant functors $E_j$ preserve restrictions, zeros, ranges,
binary restriction products and coproducts, and finite domain joins.
For $i=1,2$, the displayed Turing morphism
$\tau:V\times V\rightharpoonup V$ of \cite{ThreePoint} is sent
to a Turing morphism on $\N^2$ in $\mathcal K$.
\end{proposition}
\begin{proof}
Component evaluation preserves identities, monomorphisms, and
pullbacks in the total categories. Consequently \eqref{eq:endpoint-functor}
is independent of the representative span and preserves partial-map
composition. The stated operations are componentwise domain
restrictions, empty sets, images, products, tagged unions, and finite
unions, respectively.

Fix a recursive tuple coding and an acceptable universal partial
recursive function $\varphi_p$. At either endpoint the chosen
Turing morphism has the formula
\[
 (E_i\tau)((k,p),(\ell,x))
 =\operatorname{decode}_2\varphi_p(\operatorname{encode}_2(\ell,x)).
\]
For every partial recursive
$F:Z\times\N^2\rightharpoonup\N^2$ with c.e.\ carrier $Z$,
specialization gives a total-on-$Z$ recursive index map $s$ for
the functions $F(z,-)$, including their domain tests.
The total parameter map $z\mapsto(0,s(z))$ proves the required
universality.
\end{proof}

Let $H_i$ have singleton component at endpoint $i$ and empty
components elsewhere. An \emph{endpoint probe} of $Y$ is a total
arrow $p:H_i\to Y$. Every such arrow is uniquely
the point $p_y$ determined by an element $y\in Y_i$. A partial
arrow from $H_i$ is either total or zero, since $H_i$ has just
two restriction domains. Hence postcomposition defines a partial
function on total probes, defining $\Phi_i:\mathcal A\to\Pfn$ by
\[
 \Phi_i(Y)=\operatorname{Tot}(\mathcal A)(H_i,Y),\qquad
 \Phi_i(f)(p)=fp\quad\hbox{when }fp\hbox{ is total}.
\]
If $fp$ is not total, it is zero, and every further composite
remains zero. If it is total, further postcomposition has exactly
the domain specified by the next partial function. This proves
$\Phi_i(gf)=\Phi_i(g)\Phi_i(f)$; identities are immediate.
The bijections $p_y\leftrightarrow y$ identify $\Phi_i$ with the
underlying partial-set functor of $E_i$. Giving each probe $p_y$
the numerical name $y$ recovers its effective presentation in
$\mathcal K$.
For a domain $e$ on $Y$, its observed carrier is therefore
\begin{equation}\label{eq:probe-domain}
 |e|_i:=\{y\in Y_i:ep_y=p_y\}
       =\{y\in Y_i:E_i(e)(y)=y\}.
\end{equation}
This subset is enumerated by the component domain of $e$;
$E_i(e)$ is its partial identity.

Endpoint universality is specific to the chosen components.
For example, the numerical total function
$F(k,z,a,b)=(0,0,z,0)$ on $V_0$ cannot equal
$\tau_0(p,-)$ for a fixed central program $p$.
Inputs differing only in $z$ have the same first endpoint input
$(k,a)$, so that endpoint program forces the same endpoint output.
The first endpoint of $F(k,z,a,b)$ is $(0,z)$, a contradiction.
Thus $E_0$ preserves the listed categorical structure but does
not preserve this Turing morphism's universality.

\subsection{Proof observations and extensions}

Write $1$ for the terminal diagram with singleton components.
A \emph{global point} of $V$ is a total arrow $1\to V$.
For $j=0,1$, let $q_j$ be the range of the global point
with central value $(j,0,0,0)$ and both endpoint values $(j,0)$.
Put
\[
 d=\tau\Delta_V,\qquad a=d^*q_0,\qquad b=d^*q_1.
\]
The two domains are disjoint. For an endpoint $i\in\{1,2\}$,
write $d_i=E_i(d)$ and let $A_i,B_i\subseteq\N^2$ be the
carriers of $E_i(a),E_i(b)$. By Proposition~\ref{prop:endpoint},
\[
 A_i=\{x:d_i(x)=(0,0)\},\qquad
 B_i=\{x:d_i(x)=(1,0)\}.
\]
These sets are recursively inseparable. Indeed, if a recursive
$C\subseteq\N^2$ separated them, the total recursive function
$h$ taking value $(1,0)$ on $C$ and $(0,0)$ off $C$ would have
a program parameter $p$ with $(E_i\tau)(p,x)=h(x)$.
Then $d_i(p)=h(p)$ places $p$ in $B_i$ when $p\in C$
and in $A_i$ when $p\notin C$, contradicting separation.
A complemented separator of $a,b$ in $\mathcal A$ would be
sent to a recursive separator of $A_i,B_i$.

Use the consistent classical presentation constructed in
\cite[Theorem 6.3]{ThreePoint}. Define the covariant polynomial
$P:\mathcal E\to\mathcal E$ on the total diagram category by
$P(Z)=V+Z+(Z\times Z)$ and
$P(f)=1_V+f+(f\times f)$ for total arrows $f$.
Its initial algebra $\iota:P(S)\to S$ is constructed from finite
formula trees. Initiality means that every total algebra
$\alpha:P(Z)\to Z$ admits a unique total arrow $u:S\to Z$
with $u\iota=\alpha P(u)$.
The three summands give constructors for atoms, negation, and implication.
Write $\theta:V\to S$ for the total atom map
and $\nu:S\to S$ for the total negation constructor.
Let $\mathrm{Prf}$ be the object of finite proofs for the rules below
and $c:\mathrm{Prf}\to S$ its total conclusion map.
The theorem domain is $t_0=\ran(c)$.
Writing $\alpha\Rightarrow\beta$ for the implication constructor
applied to $(\alpha,\beta)$, the logical axiom schemes are
\[
 \alpha\Rightarrow(\beta\Rightarrow\alpha),\qquad
 (\alpha\Rightarrow(\beta\Rightarrow\gamma))
   \Rightarrow((\alpha\Rightarrow\beta)\Rightarrow
                         (\alpha\Rightarrow\gamma)),\qquad
 (\nu\beta\Rightarrow\nu\alpha)\Rightarrow(\alpha\Rightarrow\beta).
\]
Their parameters range over $S$. Modus ponens concludes $\beta$
from $\alpha$ and $\alpha\Rightarrow\beta$.
The additional atomic axioms are indexed by the split objects
of $a$ positively and $b$ negatively.
The one-node proofs and the consistency theorem give
\begin{equation}\label{eq:atomic-inequalities}
 a\le\theta^*t_0,\qquad
 b\le\theta^*\nu^*t_0,\qquad t_0\,\nu^*t_0=0.
\end{equation}
A \emph{computational extension} is a domain $t$ on $S$ that
contains $t_0$ and is closed under those logical rules.
It is consistent when $t\,\nu^*t=0$, and complete when
$t\vee\nu^*t=1_S$. The inverse image $\nu^*t$ records formulas
whose negations belong to $t$.

\begin{theorem}\label{thm:cosheaf-observation}
Fix either endpoint $i$. For every consistent computational
extension $t$ of the preceding presentation, define
\[
 P_{i,t}=(E_i\theta)^{-1}(|t|_i),\qquad
 N_{i,t}=(E_i\theta)^{-1}(|\nu^*t|_i).
\]
These are disjoint c.e.\ subsets of $\N^2$ satisfying
$A_i\subseteq P_{i,t}$ and $B_i\subseteq N_{i,t}$.
They are recursively inseparable. The assignment
\[
 t\longmapsto(\N^2,P_{i,t},N_{i,t})
\]
with identity maps on $\N^2$ for extension inclusions is a
covariant functor from the poset of consistent computational
extensions of $t_0$ to $\operatorname{Pair}(\mathrm{Rec})$.
Every consistent computational extension is incomplete.
\end{theorem}
\begin{proof}
The atom map $E_i\theta:\N^2\to S_i$ is total recursive.
Proposition~\ref{prop:endpoint} makes the two observed theorem
domains c.e.; total inverse image therefore makes $P_{i,t}$
and $N_{i,t}$ c.e. Consistency and preservation of meets and
zero make them disjoint. Applying $E_i$ to
\eqref{eq:atomic-inequalities} and using $t_0\le t$ gives
the two inclusions. A recursive separator would consequently
separate $A_i,B_i$, which is impossible.

If $t\le t'$, monotonicity gives inclusions of both observed
sets, so the identity of $\N^2$ is the required pair morphism.
The functor laws follow from the identity maps. If $t$ were
complete, preservation of joins would give
$P_{i,t}\cup N_{i,t}=\N^2$. Dovetailing their enumerations would
then decide $P_{i,t}$, contradicting inseparability.
\end{proof}

The construction uses the computational domains of
$\mathcal A$. Its numerical observations are obtained by an
actual functor on partial arrows and a total atom map.
The restriction $E_i|_{\mathscr R_X}$ starts in the nonnumerical
recursion category; the splitting accommodates its syntax and
proof objects. Extending this method to other recursion categories
requires an effective observation functor that preserves a specified
inseparable pair and the atomic proof inclusions.

\begin{thebibliography}{99}
\bibitem{Bernardi}
C.~Bernardi, \emph{Precomplete equivalence relations in dominical
categories}, Proc. Amer. Math. Soc. \textbf{98} (1986), 124--134.
\bibitem{CockettLack}
J.~R.~B. Cockett and S.~Lack, \emph{Restriction categories. I.
Categories of partial maps}, Theoret. Comput. Sci. \textbf{270}
(2002), 223--259.
\bibitem{DPH}
R.~A. Di Paola and A.~Heller, \emph{Dominical categories:
recursion theory without elements}, J. Symbolic Logic
\textbf{52} (1987), 594--635.
\bibitem{DPM}
R.~A. Di Paola and F.~Montagna, \emph{Some properties of the syntactic
$P$-recursion categories generated by consistent, recursively enumerable
extensions of Peano arithmetic}, J. Symbolic Logic \textbf{56} (1991),
643--660.
\bibitem{Heller}
A.~Heller, \emph{An existence theorem for recursion categories},
J. Symbolic Logic \textbf{55} (1990), 1252--1268.
\bibitem{Lengyel2004}
F.~Lengyel, \emph{More existence theorems for recursion categories},
Ann. Pure Appl. Logic \textbf{125} (2004), 1--41.
\bibitem{Corrections}
F.~Lengyel, \emph{Corrections to More existence theorems for recursion
categories}, version 10, 20 October 2026,
\href{https://github.com/flengyel/errata/blob/5555e4fa1da7b8e3a0836dd1fbf1ca385cacd03b/lengyel2004/lengyel2004_repairs_v10.tex}{source}.
\bibitem{Composition}
F.~Lengyel, \emph{Pointwise provable equality and the failure of
composition}, 2026,
\href{https://arxiv.org/abs/2609.25556}{arXiv:2609.25556}.
\bibitem{LOC}
F.~Lengyel, \emph{Locally connected recursion categories}, manuscript,
revised 10 October 2026, \texttt{LOCOct10f.tex}.
\bibitem{ThreePoint}
F.~Lengyel, \emph{Range completion and computational cosheaves on a
three-point space}, 29 October 2026,
\href{https://github.com/flengyel/CategoricalEssentialIncompleteness/blob/3be65406b0bb6d3109a41fd3f517ca62fcc5ed8f/manuscript/range_completion_three_point_cosheaves.tex}{source}.
\bibitem{Montagna}
F.~Montagna, \emph{``Pathologies'' in two syntactic categories of partial
maps}, Notre Dame J. Formal Logic \textbf{30} (1989), 105--116.
\bibitem{RosoliniThesis}
G.~Rosolini, \emph{Continuity and effectiveness in topoi}, D.Phil.
thesis, University of Oxford, 1986.
\bibitem{Rosolini}
G.~Rosolini, \emph{Representation theorems for $P$-categories},
in F.~Borceux (ed.), \emph{Categorical Algebra and its Applications},
Lecture Notes in Mathematics \textbf{1348}, Springer, Berlin, 1988,
307--315.
\end{thebibliography}
\end{document}
